\section{A cancellation-preserving Gaussian derivative certificate}
\label{sec:joint-gaussian-norm}

I refine the scalar certificate without changing the classical financial
model or the Gaussian-tilt coefficients. The signed-Stirling relation between
price and log-price derivatives has an explicit financial antecedent in
\citet{carr2001}, equation (9). Conditional Gaussian polynomial basket
pricing is also established \citep{olivaresalvarez2016}. The object I check
here is the norm of the combined derivative, followed by its concrete effect
on the finite financial remainder. I distinguish it from a sum of separately
bounded derivatives.

Let $\nu$ be a finite nonnegative Borel measure on real log thresholds, with
mass $H>0$, and let $\lambda>0$. Define
\[
 P_\nu(m)=\int_{\mathbb R}\Phi((m-x)/\lambda)\,\nu(dx),\qquad
 g_m(t)=P_\nu(m+\log t),\quad t>0.
\]
For the capped call with strike and reference spot $K$, the measure is
$\nu(dx)=K e^x\mathbf1_{[0,\log(1+H/K)]}(x)\,dx$, of mass $H$.
A cash digital corresponds to a point mass. Positive finite mixtures of
such threshold payoffs are therefore included. This statement does not
assume an unbounded call payoff, negative basket weights or a signed
threshold measure.

For integer $r\geq1$, let
\begin{equation}
 A_{r,\lambda}(z)=\sum_{\ell=1}^r s(r,\ell)\lambda^{-\ell}
              (-1)^{\ell-1}\operatorname{He}_{\ell-1}(z),\qquad
 D_r(\lambda)=\sup_{z\in\mathbb R}|A_{r,\lambda}(z)\phi(z)|.
 \label{eq:joint-gaussian-norm}
\end{equation}
The Hermite identity
$\phi^{(j)}(z)=(-1)^j\operatorname{He}_j(z)\phi(z)$ follows by induction
from $\phi'=-z\phi$ and the polynomial recurrence. Every polynomial times
$\phi$ is bounded and tends to zero at both infinities. For each finite
derivative order, dominated differentiation under $\nu$ is consequently
valid because its total mass is finite. Applying the already derived
signed-Stirling identity gives
\begin{equation}
 g_m^{(r)}(t)=t^{-r}\int A_{r,\lambda}((m+\log t-x)/\lambda)
                    \phi((m+\log t-x)/\lambda)\,\nu(dx).
 \label{eq:joint-gaussian-integral-derivative}
\end{equation}

\begin{theorem}[A sharp threshold-measure derivative bound]
\label{thm:joint-threshold-norm}
For $t>0$, $|g_m^{(r)}(t)|\leq H D_r(\lambda)t^{-r}$.
The constant $H D_r(\lambda)$ is the smallest uniform constant over all
nonnegative threshold measures of mass $H$ and all $m$ and $t$.
For the Gaussian basket with $T\geq1$, the same coefficients as in
Theorem~\ref{thm:finite-gaussian-price-certificate} satisfy
\begin{equation}
 \left|V-\sum_{j=0}^k c_j\right|
 \leq {H D_{k+1}(\lambda)\over(k+1)!}
                      \mathbb E[(T-1)^{k+1}].
 \label{eq:joint-norm-price-remainder}
\end{equation}
\end{theorem}
\begin{proof}
Taking absolute values in
equation~\eqref{eq:joint-gaussian-integral-derivative} and using the
nonnegative measure gives the upper bound. The nonzero polynomial
$A_{r,\lambda}$ has leading coefficient
$(-1)^{r-1}\lambda^{-r}$, because $s(r,r)=1$ and the Hermite polynomial
is monic. Its Gaussian-weighted absolute maximum is positive and attained:
continuity and decay place every positive maximum inside a compact interval.
Choose a maximiser $z_*$, $t=1$, any threshold $x_*$, $m=x_*+\lambda z_*$
and $\nu=H\delta_{x_*}$. The integral then has absolute value exactly
$H D_r$. No smaller uniform constant can work. This attainment is over
the stated measure class; the particular capped-call measure need not
attain it. Taylor's integral remainder on $[1,T]$ is bounded using
$t^{-r}\leq1$, and its expectation is finite under the Gaussian model.
This proves the last assertion without altering a coefficient or assuming
convergence of the infinite expansion.
\end{proof}

\subsection{A complete rational root and outward arithmetic certificate}

Put $Q_{r,\lambda}=A'_{r,\lambda}-zA_{r,\lambda}$. Its degree is $r$.
At every nonzero absolute maximum of $A\phi$, differentiation gives
$(A\phi)'=Q\phi=0$. Zeros of $A$ have zero absolute value, and the
limits at infinity are zero. Thus it suffices to evaluate every real
root of $Q$. When $\lambda$ is rational, every coefficient is rational.
I remove multiplicities exactly by replacing $Q$ with
$q=Q/\gcd(Q,Q')$, with a positive-leading monic normalisation.
Removing multiplicities preserves the set of distinct roots.

For completeness I spell out the root-counting certificate. Start with
$q_0=q$, $q_1=q'$ and let $q_{j+1}$ be the negative remainder in division
of $q_{j-1}$ by $q_j$. Positive rescaling of each new polynomial leaves
all signs unchanged. Square-freeness makes the final polynomial a nonzero
constant. Let $V(x)$ be the number of sign variations, omitting zeros.
Consecutive members cannot vanish together: the division recurrence would
propagate that zero to the final constant. At a zero of an internal member,
its neighbours have opposite signs by the same recurrence, so their two
adjacent comparisons always contribute one variation. Passing that zero
does not change $V$. At a zero of $q_0$, $q_1=q'$ is nonzero, and $q_0$
changes from the opposite sign to the same sign as $q_1$; $V$ decreases by
one. These are all possible changes. Therefore $V(a)-V(b)$ counts the
distinct roots in $(a,b)$ when the endpoints are not roots.

For a degree-$d$ polynomial with leading coefficient $a_d$, every root
has modulus at most $1+\max_{j<d}|a_j/a_d|$. To check this bound, outside
that radius the absolute sum of the lower terms is at most
$|a_d|M(|z|^d-1)/(|z|-1)<|a_d z^d|$, so the leading term cannot cancel.
I use a strictly larger integer bound. Exact rational subdivision retains
only intervals with a positive variation difference, until each contains
one root and has width at most $2^{-120}$. If a trial midpoint is a root,
I select a nonroot interior rational point. A polynomial of degree $d$
cannot vanish at all $d+1$ proposed alternatives. I take these alternatives
at fractions $j/(d+2)$ of the interval, $1\leq j\leq d+1$.
Every resulting child has width at most $(d+1)/(d+2)$ times its parent
width; an ordinary midpoint split has factor $1/2$. The contraction is
uniformly strictly below one at a fixed degree. Every split preserves
the sum of root counts. Separate roots have positive separation, so
contraction first isolates them and then reaches the required width.
This proves termination, including intervals whose midpoint is a root.

For each root interval $[a,b]$, I construct an Arb ball containing both
exact rational endpoints and evaluate $|A(z)\phi(z)|$ at 384 bits on that
ball. Its outward upper endpoint bounds the value at the root. The largest
of these endpoints is a global norm upper certificate. Function evaluations
at interval midpoints provide lower bounds on the global supremum, even
when those midpoints are not roots. The independent implementation uses
Flint rational polynomials reconstructed from falling factorials and closed
Hermite coefficients, counts roots from signs at both infinities, checks
every interval and verifies every Arb box against its exact endpoints.
Fresh 640-bit evaluations provide an additional envelope check.

I also certify the separate norms
$\mathcal M_j=\sup_z|\operatorname{He}_j(z)\phi(z)|$. The three distinct comparisons
are
\[
 D_r\leq E_r^{\rm tri}:=
       \sum_{\ell=1}^r|s(r,\ell)|\lambda^{-\ell}\mathcal M_{\ell-1}
       \leq C_r.
\]
The first inequality is the triangle inequality; the second follows from
the original Fourier bounds. This intermediate opponent distinguishes
cancellation from improving the individual derivative envelopes. I retain
all three certified bounds, not just their largest numerical difference.

\subsection{Matched finite pricing consequences}

All 27 coupled constants, nine individual Hermite norms and 180 real
stationary points pass the independent audit. Relative to the original
triangle constants, the coupled upper bounds range from equality at order
one to a ratio of approximately $0.088959$. Against the sharper individual
norm triangle, the smallest ratio is approximately $0.094075$.
These are derivative constants, not pricing errors or running times.

I reuse the retained exact coefficient intervals, recompute their moment
geometry and tighten all nine degree intervals. The same 25 contracts
certify; the two weak-smoothing controls remain unresolved by this family
of Taylor bounds. Three first certifying degrees decrease:
the four-asset $\lambda=0.2$, $H=10$ case from two to one;
the eight-asset $\lambda=0.05$, $H=1$ case from two to one; and its
$H=25$ case from four to three. I then execute the stopping hierarchy with
fresh geometry banks. It uses 793 scalar kernels and 760 geometric
exponentials, compared with 1,169 and 1,552 under the preceding bound.
The constant-generation run takes approximately $4.52$ seconds; the
executed hierarchy takes approximately $0.156$ seconds. The generation
cost is additional shared setup. A single cold run is therefore not
established to be faster in total host time, and the counts do not
constitute superiority over the best published pricer.

The raw-remainder divergence proposition remains a statement about its
original $C_r$ bound. I do not silently transfer it to $D_r$. The sharp
operator statement is likewise not sharpness of the expected financial
remainder. Its role in the revised quantum design is to supply a rigorous
normalisation for a signed density-derivative oracle.
